\documentclass[10pt,a4paper]{amsart}
\usepackage[margin=19mm]{geometry}
\usepackage{amsmath,amssymb,mathtools}
\usepackage[scr=rsfs,cal=euler]{mathalfa}
\usepackage{xfrac}
\usepackage[unicode,hidelinks,bookmarks=false]{hyperref}
\newcommand{\ie}{i.e.\ }
\newcommand{\cf}{cf.\ }
\newcommand{\resp}{respectively\ }
\newcommand{\Subsection}{\subsection}
\providecommand{\nobreakdash}{\nobreak\hbox{-}}
\setlength{\emergencystretch}{2em}
\multlinegap0pt
\usepackage{mathtools}
\usepackage[shortlabels]{enumitem}
\usepackage{float}
\usepackage{tcolorbox}
\usepackage{amssymb}
\newtheorem{theorem}{Theorem}
\title{A lifting-projection numerical framework for Boltzmann and Landau equations}
\author{Kun Huang}
\date{Source: arXiv:2605.03273v1 (CC BY 4.0)}
\begin{document}
\maketitle

\noindent\emph{Source and licence.} A lifting-projection numerical framework for Boltzmann and Landau equations. Version arXiv:2605.03273v1 (CC BY 4.0), distributed by arXiv under the Creative Commons Attribution 4.0 International licence, http://creativecommons.org/licenses/by/4.0/, as recorded in the arXiv licence metadata for this exact version and shown on the abstract page https://arxiv.org/abs/2605.03273v1. Full licence notice: \texttt{licenses/arxiv-2605.03273-CC-BY-4.0.txt}.\par\smallskip

\noindent\textit{Editorial scope.} Counted unit: the article's theorem thm:entropy (source0.tex lines 420--438). The article's complete original proof of it is delivered with it (source0.tex lines 439--489, one \texttt{proof} environment of 51 lines). No mathematics is written, repaired or re-derived: the statement and the proof are the article's own lines, in the article's own notation, and no retained line is substituted or re-flowed.  No further source-local context is needed: every cross-reference of every retained line resolves inside the two delivered spans.  Nothing was omitted from any retained span, so the package carries no omission record.  No retained line contains a citation, so the wrapper carries no bibliography and no bibliography entry.  No retained line contains a figure environment, an image-inclusion command or drawing code, so no image is delivered and the package declares no asset dependency.  Editorial formatting only, in addition: an AMS \texttt{amsart} preamble that restates the article's own declarations and macros used by the retained lines, namely the environments \texttt{theorem} on separate counters, together with the standard packages those lines need.  The retained environments print in this document as Theorem 1 in order of appearance, whereas the article prints the same items with the numbering of its own full document.  The complete native source of the article as distributed in the arXiv e-print for this exact version is bundled unchanged as \texttt{sources/199780/source0.tex}.
\begin{theorem}[H-theorem for consecutive LP flow]\label{thm:entropy}
Define the entropy $H$ of distribution functions as:
\begin{equation*}
    \begin{split}
        H(f(\mathbf{v})) \coloneqq & \int_{v} f(\mathbf{v})\log{f(\mathbf{v})},\\
        H(F(\mathbf{v},\mathbf{w})) \coloneqq & \iint_{vw} F(\mathbf{v},\mathbf{w})\log{F(\mathbf{v},\mathbf{w})}.
    \end{split}
\end{equation*}
    Let $(t, F_{K}(t))$ be the consecutive 2-particle Kac flow, and $(t, \widetilde{f}(t))$ be the consecutive LP flow. 
    
    For any $0\leq t_{1}<t_{2}<+\infty$ we have
    \begin{equation*}
            H(F_{K}(\cdot, t_{2})) \leq H(F_{K}(\cdot, t_{1})).
    \end{equation*}
    In addition, for any $0\leq n < m < \infty$ we have
    \begin{equation*}
        H(\widetilde{f}^{(m)}) \leq H(\widetilde{f}^{(n)}). 
    \end{equation*}
\end{theorem}

\begin{proof}
    Within each interval we know that the 2-particle Kac master equation satisfy the H-theorem:
    \begin{equation*}
        \partial_{t} H(F) \leq 0.
    \end{equation*}
    It suffices to show that
    \begin{equation*}
        H\left(F_{K}(n\Delta t)\right)  \leq H\left(\lim_{t\rightarrow (n\Delta t)^{-}} F_{K}(t)\right).
    \end{equation*}

    Note that
    \begin{equation*}
        F_{K}(n\Delta t) = \left( \lim_{t\rightarrow (n\Delta t)^{-}}\pi F_{K}(t)\right) \otimes \left( \lim_{t\rightarrow (n\Delta t)^{-}}\pi F_{K}(t)\right) \neq \lim_{t\rightarrow (n\Delta t)^{-}} F_{K}(t),
    \end{equation*}
    hence it remains to show that for any $F$ s.t. $F(\mathbf{v}, \mathbf{w}) = F(\mathbf{w}, \mathbf{v})$ we have
    \begin{equation*}
        H(\pi F \otimes \pi F) \leq H(F).
    \end{equation*}

    Recall that the mutual information is always non-negative:
    \begin{equation*}
        I(F) \coloneqq \iint_{vw} F \log{\frac{F}{\pi F\otimes \pi F}} \geq 0.
    \end{equation*}
    Decompose the mutual information as follows
    \begin{equation*}
    \begin{split}
\iint_{vw} F \log{\frac{F}{\pi F\otimes \pi F}} = &\iint_{vw}F\log{F} - \iint_{vw} F \log{\pi F(\mathbf{v})} - \iint_{vw} F \log{\pi F(\mathbf{w})}\\
=& \iint_{vw}F\log{F} - \int_{v} \left(\int_{w}F\right) \log{\pi F(\mathbf{v})} - \int_{w} \left(\int_{v}F\right) \log{\pi F(\mathbf{w})}\\
=& \iint_{vw}F\log{F} - 2\int_{v} (\pi F) \log{(\pi F)}
    \end{split}
    \end{equation*}

    Now we can see that 
    \begin{equation*}
        H(F) \geq 2 H(\pi F) = H(\pi F \otimes \pi F).
    \end{equation*}
    Consequently
    \begin{equation*}
        H\left(F_{K}(n\Delta t)\right) = H\left(\left( \lim_{t\rightarrow (n\Delta t)^{-}}\pi F_{K}(t)\right) \otimes \left( \lim_{t\rightarrow (n\Delta t)^{-}}\pi F_{K}(t)\right)\right) \leq H\left(\lim_{t\rightarrow (n\Delta t)^{-}} F_{K}(t)\right).
    \end{equation*}

    Recall that $\widetilde{f}^{(n)} = \lim_{t\rightarrow (n\Delta t)^{-}}\pi F_{K}(t)$ and $2 H(\pi F) = H(\pi F \otimes \pi F)$, therefore the previous conclusion
    \begin{equation*}
        H\left(\left( \lim_{t\rightarrow (m\Delta t)^{-}}\pi F_{K}(t)\right) \otimes \left( \lim_{t\rightarrow (m\Delta t)^{-}}\pi F_{K}(t)\right)\right) \leq  H\left(\left( \lim_{t\rightarrow (n\Delta t)^{-}}\pi F_{K}(t)\right) \otimes \left( \lim_{t\rightarrow (n\Delta t)^{-}}\pi F_{K}(t)\right)\right)
    \end{equation*}
    implies that
    \begin{equation*}
        H(\widetilde{f}^{(m)}) \leq H(\widetilde{f}^{(n)}).
    \end{equation*}
    
    \end{proof}

\end{document}
